\begin{theorem}[the actual smooth monodromy is trivial]\label{push:src:C.3.24}
For the actual retained $S^6$ family and a positive generator
of its punctured $w$-parameter disc, the oriented smooth
monodromy satisfies
\[
 [\Phi]=0\quad\hbox{in }\pi_0\operatorname{Diff}^+(S^6)
                       \simeq\Theta_7\simeq\Z/28.
\]
Its mapping torus is consequently oriented-diffeomorphic to
$S^1\times S^6$. Every finite cyclic base change has trivial
smooth monodromy as well. By \ref{push:src:C.4.1}'s marked family comparison,
the same conclusion applies to the admissible sphere twist
data covered by that comparison.
\end{theorem}
\begin{proof}
Let $\widehat\psi$ be the actual smooth equivariant lift of
the annular orbit map in \ref{push:src:C.3.10}.
The actual cap calculation is an integer calculation.
\ref{push:src:C.3.18} gives $p_{\mathrm{rel}}=-24$ and \ref{push:src:C.3.21} gives
$s_\delta=-2$. The common signed comparison in \ref{push:src:C.3.22} yields
\[
 J_\psi=\epsilon_{\rm cap}
          \frac{-24-12(-2)}{24}=0\quad\hbox{in }\Z.
\]
By \ref{push:src:C.3.14} and \ref{push:src:C.3.23}, the unframed orbit-knot integer of the
actual annular lift is therefore zero. Neither the overall cap
orientation sign nor the Wall--Davis generator sign can change
this conclusion.

We verify that this is the sixth power at the equivariant
integer level. The family carries the fixed positive
$\delta$ circle, and all its closing maps commute with that
circle. Choose a circle-invariant Ehresmann connection along
one positive $w$ loop. Concatenating it six times gives the
sixth power of the same equivariant return diffeomorphism.
Changing to the invariant transport used in \ref{push:src:C.3.10} changes
this return by an equivariant isotopy: invariant horizontal
connections form an affine space, and their interpolated
flows give that isotopy.

The closing cocycle itself introduces no hidden discrepancy.
In the notation of Paper~I, Proposition~18.5,
\[
 b_{h,c}=-\frac{h}{12}\Pi_c\widehat\gamma
             +\left(h\ell-\frac{h^2}{24}\right)e,
 \qquad \Pi_{c+h}\widehat\gamma
                   =\Pi_c\widehat\gamma+h e.
\]
Direct substitution gives
\[
 b_{h,c}+b_{k,c+h}=b_{h+k,c}.
\]
Thus six $h=2$ closings are the $h=12$ closing used for the
annular reduction. Returning from these coherent closings to
any stipulated Paper~I, Proposition~3.3 and Paper~I, Proposition~3.4 and Paper~I, Theorem~3.8 choices changes them by a scalar in the
connected group $\C^*$, which centralizes $S^1_\delta$;
this is again an equivariant isotopy. Every subsequent unit isotopy, finite-chart
transport and circle-gauge comparison in \ref{push:src:C.3.6} and \ref{push:src:C.3.10} is
equivariant. Consequently the annular lift and $\Phi^6$ have
the same integer $\mathfrak q$, not merely the same underlying
ordinary mapping class.

Additivity from \ref{push:src:C.3.23} now gives
\[
 6\mathfrak q(\Phi)=\mathfrak q(\Phi^6)
                   =\mathfrak q(\widehat\psi)=-J_\psi=0.
\]
Since this equality is in $\Z$, it implies
$\mathfrak q(\Phi)=0$. Its signed reduction modulo $28$
is the ordinary class $[\Phi]$, so that class is zero.
This removes the possible order-two ambiguity by an integer
argument; it does not divide the equation $6[\Phi]=0$ inside
$\Z/28$.

A smooth isotopy from the return map to the identity gives
the usual oriented product identification of its mapping
torus. A finite cyclic base change replaces the return map
by a positive power, whose class is again zero. The final
twist-data statement follows from the local family isomorphisms
and connected automorphism transitions already proved in \ref{push:src:C.4.1}.
\end{proof}
